\chapter{Introduction}
\section{Motivation}
Topic of this thesis is motivated, by my work in Formula Student team eForce
Prague Formula, where I worked on real-time vehicle dynamics control software.
There I got intricately exposed to problems faced by racing teams, such as
should we increase the aerodynamic downforce at the cost of higher drag? Should
we put accumulator in the middle of vehicle, to decrease moment of inertia?
Should the aerodynamics be designed for straight line drive, or for steady
state turning? How to use limited energy from accumulator more effectively, how
faster can we go if we made better power limiting algorithm? How large the
accumulator should be? Which torque allocation algorithm should be used? How to
design the suspension kinematics to maximize the tire utilization? Is rear
wheel steering worth the increased weight? All of these problems can be
analyzed by using purpose built tool, good engineering practices or using pen
and paper with heavily simplified model. Formula student season is 10 months
for development and construction of the vehicle, 2 months for competition. That
puts enormous pressure on deciding on vehicle concept fast. In my opinion what
formula student teams need is a tool to quickly evaluate different car
concepts. As my bachelor thesis I created a prototype of such tool, goals of
this thesis come from computational and user experience form my previous tool.

\section{Goals}
\begin{itemize}
	\item Primary goal of this thesis is to investigate available methods for numerical
	      optimal control with respect to their suitability for planning an optimal
	      trajectory of a vehicle aiming to achieve the minimum lap time on a racing
	      circuit. To do this for various vehicle concepts, thesis has to present a
	      vehicle modelling framework.
	\item Benchmark performance of available numerical methods
	\item Thesis should present a tool for convenient vehicle modelling, and vehicle
	      concept evaluation.
	\item Support vehicle parameter optimization, using any vehicle parameter as control
	      input and sensitivity analysis of vehicle parameters.
	\item Goal is not precise vehicle modeling and validation, goal is to create a
	      framework to perform vehicle concept analyses easily.
\end{itemize}

%Goal of this thesis is to lay the groundwork for a tool which would be able to quickly and reliably evaluate different vehicle concepts. 
%Thesis should presenta framework for simulating multiple complexities of vehicle model. 
%like in bachelor thesis where i preformeda nd vladiated analysis on vehicle components, 

%Each parameter can be static, desgin varaible, or control variable.

%For fast and sufficiently accurate vehicle simulation it is crucial to explore and compare
%various methods of transcribing the optimal control problem

% such as global collocation using a single polynomial over the whole horizon, multi-segment collocation that splits the horizon into smaller intervals each with its own polynomial, and Runge Kutta methods.

%Goal is not to present results of specific car model or recommendation on car
%design, therefore validation and verification were not done in this thesis. i
%ahev already done them in bachelor thesis
%\cite{ciglanskyAnalysisVehicleComponents}

\section{Structure of thesis}
First part of thesis with vehicle modelling in the context of optimal control.
A universal model structure inspired by \cite{brown_engineering_2006} is
presented. Then the theory of optimal control idem spat ja neviem co pisat uz.
\begin{enumerate}
	\item ODE models of cars
	\item Solving methods: time marching forward, backward pass, NLP,
	\item Collocation (p, h, hp refinement), integral + differential form (Gauss-Radau,
	      Legendre, Lobatto)
	\item Implementation overview, Algorithmic differentiation, JUMP/Examples/InfiniteOpt
	      + Julia, Python/MATLAB, CasADi, interior point solvers
	\item Actual implementation and results
	\item Conclusion
\end{enumerate}

\section{Literature overview/ survey of laptime simulators}


Lap-time prediction dates back to the 1950s. Computers became common in the
1970s and quasi-steady-state (QSS) methods became the standard approach for
MLTS.\@ QSS is still what most race engineers mean by ``lap-time simulation''.
In QSS the vehicle is a point mass on a fixed racing line, with lateral and
longitudinal accelerations bounded by a steady-state performance
envelope~\cite{Roland1971}. Transient effects such as yaw acceleration are
ignored. MLTS was first formulated as an optimal control problem in the late
1980s~\cite{Hatwal1986SomeIS}. A survey of current methods is given
in~\cite{massaroMinimumlaptimeOptimisationSimulation2021}. The work of Perantoni
and Limebeer~\cite{perantoniOptimalControlFormula2014} is the main reference:
it uses a track-fixed coordinate system, switches the independent variable from
time to travelled distance, and applies the formulation to a twin-track
Formula 1 model with some vehicle parameters as decision variables. Parts of
this thesis build on it. MLTS is still an active research area because better
predictions give measurable setup gains, but higher fidelity models cost more
solver time. The authors of~\cite{ComputationallyEfficientFramework} use
sequential convex programming to cut compute cost without losing accuracy.

A very popular laptime simulator is OptimumLap from OptimumG
\cite{OptimumLapOptimumG}, It is a free commercial tool advertised to achieve
10\% accuracy with real vehicles. It uses Quasi Steady State methods (QSS) This
is the standard tool for laptime simulations. Open source implementation of the
same method \cite{mc12027Mc12027OpenLAPLapTimeSimulator2026}, My implementation of this method in chapter ??. QSS methods are very fast, but the models for which they can be used is equally very simple.

Popular open source C++ implementation of laptime simulator using nonlinear
programming and automatic differentiaoton\cite{manzaneroJuanmanzaneroFastestlap2026}. It uses IPOPT and
CppAD\cite{CoinorCppAD2026}. Vehicle model is 6DOF twintrack. Other implementations are
work from austria or germany\cite{Heilmeier2019} class project from Georgia institute of technology \cite{2026ClassProject}
